\documentclass[handout]{beamer}

\mode<presentation>
{
  \usetheme{default}      % or try Darmstadt, Madrid, Warsaw, ...
  \usecolortheme{default} % or try albatross, beaver, crane, ...
  \usefonttheme{default}  % or try serif, structurebold, ...
  \setbeamertemplate{navigation symbols}{}
  \setbeamertemplate{footline}[frame number]
  \setbeamertemplate{caption}[numbered]
}

\usepackage[english]{babel}
\usepackage[utf8x]{inputenc}

\usepackage{graphicx}
\usepackage{subcaption}
\usepackage{hyperref}

\title{Math Camp 2020: Programming (part 1)}
\author{Frank Pinter}
\date{20 August 2020}

\AtBeginSection[]
{
  \begin{frame}
    \frametitle{Table of Contents}
    \tableofcontents[currentsection]
  \end{frame}
}

\begin{document}

\begin{frame}
  \titlepage
\end{frame}

\begin{frame}{Outline}
  \tableofcontents
\end{frame}

\section*{Introduction}

\begin{frame}{Structure of today's session}
\begin{enumerate}
    \item General programming advice ($\approx 30$ min)\itemsep=2em
    \item Julia walkthrough ($\approx 60$ min)
    \begin{itemize}
        \item Feel free to ask questions throughout!
    \end{itemize}
\end{enumerate}
\vspace{0.5in}
All materials from today's presentation are on GitHub at \\
\url{https://github.com/fpinter/math-camp-coding}
\end{frame}

\section{Basic principles}

\begin{frame}{Basic principles}
    \begin{enumerate}
        \item Learn by doing (practice!)\itemsep=2em
        \item Always keep your future self in mind
        \item Your time is valuable
        \item Don't reinvent the wheel
    \end{enumerate}
\end{frame}

\begin{frame}{Learn by doing}
    \begin{itemize}
        \item Especially early in grad school, treat learning about programming as an investment\itemsep=2em
        \item Practice new skills as often as you can
        \item Programming needs frequent reinforcement
    \end{itemize}
\end{frame}

\begin{frame}{Always keep your future self in mind}
    \begin{itemize}
        \item When you return to a project later on, you should be able to:\itemsep=1.5em
        \begin{itemize}
            \item Figure out what's going on relatively quickly
            \item Not screw things up
        \end{itemize}
        \item Write clear documentation and keep it updated (don't rely on your memory)
        \item Clearly written code $\gg$ over-commenting
        \begin{itemize}
            \item Use good variable names
            \item Use good function names
            \item Use functions to simplify things
        \end{itemize}
        \item Write comments with a specific audience in mind
        \begin{itemize}
            \item Typically your future self and your collaborators
        \end{itemize}
    \end{itemize}
\end{frame}

\begin{frame}{Your time is valuable}
    \begin{itemize}
        \item Your time is more valuable than the computer's time\itemsep=2em
        \item Prioritize organization, readability, and clarity over fast runtime
        \item Resist the temptation to focus on runtime too early
        \begin{itemize}
            \item Wait until you've checked for accuracy
            \item Wait until you have a clear idea which parts are actually critical
        \end{itemize}
        \item Runtime is important in parts of the code you'll be running many times as part of your usual workflow (and those parts only)
    \end{itemize}
\end{frame}

\begin{frame}{Your time is valuable}
\begin{quotation}
  Programmers waste enormous amounts of time thinking about, or worrying about, the speed of noncritical parts of their programs, and these attempts at efficiency actually have a strong negative impact when debugging and maintenance are considered. We \textbf{should} forget about small efficiencies, say about 97\% of the time: premature optimization is the root of all evil. \\

  Yet we should not pass up our opportunities in that critical 3\%.
\end{quotation}
--\href{https://doi.org/10.1145/356635.356640}{Donald Knuth (1974)}
\end{frame}

\begin{frame}{Don't reinvent the wheel}
    \begin{itemize}
        \item There might be ways to solve your problem you hadn't thought of\itemsep=2em
        \item Talk to your cohort, talk to people you know, ask questions online
        \item Stay up to date on tools and the technical community
        \item If something feels like a common problem, spend time looking for a common solution
        \begin{itemize}
            \item Poll: are you familiar with regular expressions?
        \end{itemize}
    \end{itemize}
\end{frame}

\section{More specific advice}

\begin{frame}{We are more like software engineers than we think}
\includegraphics[width=\textwidth]{files/tweet.pdf}

(\href{https://twitter.com/korenmiklos/status/1294363828646543362?s=20}{source})
\end{frame}

\begin{frame}{When writing code}
    \begin{itemize}
        \item Don't repeat yourself\itemsep=2em
        \begin{itemize}
            \item Don't copy and paste; write functions
        \end{itemize}
        \item<2-> Write (and save) formal tests
        \begin{itemize}
            \item Write tests for functions \textit{when you write the functions}
            \item Write checks your data should pass
            \item Tip: accumulate a bank of test cases over time
        \end{itemize}
        \item<3-> Know and use the idioms of your language
        \begin{itemize}
            \item Know \textit{why} your code works the way it does
            \item Know the common gotchas
            \item Nothing should be magic!
        \end{itemize}
        \item<4-> Understand all unexpected results
        \begin{itemize}
            \item Check all your results for anything unusual (smell test)
            \item Learn how to read your language's error messages
        \end{itemize}
    \end{itemize}
\end{frame}

\begin{frame}{When organizing your project}
    \begin{itemize}
        \item Split your code into steps, with a clear order\itemsep=2em
        \begin{itemize}
            \item Having a master script is strongly recommended
            \item You should be able to clear all your outputs/intermediate files and run the master script
        \end{itemize}
        \item<2-> Aim for full reproducibility, including a detailed readme file instructing a replicator \textit{exactly what to do}
        \begin{itemize}
            \item Keep this file continuously updated as you work
            \item Fewer steps for the replicator = better
        \end{itemize}
        \item<3-> Don't write critical parts of your code under time pressure
        \begin{itemize}
            \item If you do, go back and clean it up later
        \end{itemize}
        \item<4-> Use version control to track changes over time
    \end{itemize}
\end{frame}

\begin{frame}{Note on choosing programming languages}
    \begin{itemize}
        \item Tired: wars between programming languages on Twitter\itemsep=2em
        \item Wired: using the right language for the task at hand
        \item There's no rule saying you have to use the same language for everything (even within a project)
        \item<2-> Questions to ask yourself:
        \begin{itemize}
            \item What do your coauthors use?
            \item In what language do you work most efficiently?
            \item What functionality do you need?
            \begin{itemize}
                \item e.g., data cleaning, web scraping, heavy computation
            \end{itemize}
        \end{itemize}
    \end{itemize}
\end{frame}

\section{Specific advice: first year}

\begin{frame}{First year is different from research}
\begin{itemize}
    \item The work you do in first year is very different from the work you'll do later on, regardless of field\itemsep=2em
    \item First year teaches you how methods work conceptually, not how to use them in research
    \item<2-> Advice: treat this as an opportunity to learn by doing
    \begin{itemize}
        \item Develop good habits
        \item Spend some time planning your workflow and your approach
        \item Decide what skills you want to learn, and take first year as an opportunity to learn them
    \end{itemize}
\end{itemize}
\end{frame}

\begin{frame}{First year vs. research}
    \begin{tabular}{p{2.5cm}|p{3.5cm}|p{3.5cm}}
    & G1 coursework & Research/real life \\
    \hline
    Day to day work &
Numerical computation with clean or simulated data &
Mostly wrangling real-world data (unless you're a theorist) \\
\hline
Maintainability &
Submit and you're done &
Return to your code many times, sometimes years later \\
\hline
Accuracy &
Nice to have &
Be obsessive about making sure your results are correct \\
    \end{tabular}
\end{frame}

\begin{frame}{First year vs. resarch}
    \begin{tabular}{p{2cm}|p{4cm}|p{4cm}}
    & G1 coursework & Research/real life \\
    \hline
Testing &
Smell test + write formal tests for basic debugging &
Smell test + write lots of formal tests \\
\hline
Collaboration &
Discuss with group, but write code independently &
Divide tasks + perhaps do code review \\
\hline
Version control &
Nice to have, but optional &
Very important
    \end{tabular}
\end{frame}

\section{Further resources}
\begin{frame}{Further resources}
\begin{itemize}
    \item \href{https://scholar.harvard.edu/ristovska/resources}{Ljubica Ristovska's presentation}\itemsep=2em
    \item \href{https://www.sas.upenn.edu/~jesusfv/teaching.html}{Jesús Fernández-Villaverde's lecture notes}
    \item \href{https://quantecon.org/}{QuantEcon}
    \item \href{https://www.iq.harvard.edu/data-science-services/workshop-materials}{Harvard IQSS training materials}
\end{itemize}
\end{frame}

\begin{frame}{How to get help}

\begin{enumerate}
    \item Check the built-in help in the language\itemsep=2em
    \item Google
    \begin{itemize}
        \item Often the result will be a Stack Overflow answer -- these are often helpful but not always
        \item Watch out for out-of-date info (especially for Julia)
    \end{itemize}
    \item Ask someone
    \begin{itemize}
        \item Plug for the econ department Slack
    \end{itemize}
    \item Ask a question on Stack Overflow
    \begin{itemize}
        \item Stack Overflow has guidance on how to ask a good question (varies by language); read that first
    \end{itemize}
\end{enumerate}
\end{frame}

\section{Julia}

\begin{frame}{Why Julia today?}
\begin{itemize}
    \item The focus of math camp: skills you'll use in first year\itemsep=2em
    \begin{itemize}
      \item Numerical computation, matrix algebra, optimization
    \end{itemize}
    \item Julia excels at these and its matrix syntax is clean
    \item<2-> Historically the dominant language for first year PhD was Matlab
    \begin{itemize}
        \item Julia syntax is closely based on Matlab\itemsep=1em
        \item Unlike Matlab, Julia is free and open-source, with a growing community, and many of the advantages of modern languages
        \item Julia is also more efficient (especially loops, optimization, and parallelization)
        \item You can switch back to Matlab anytime if you want
    \end{itemize}
\end{itemize}
\end{frame}

\begin{frame}{Alternatives to Julia}
\begin{itemize}
    \item Matlab\itemsep=2em
    \begin{itemize}
        \item Legacy code + inertia
        \item Dynare (for macro)
        \item Lacks features of modern languages
        \item Expensive outside of academia (or use Octave)
    \end{itemize}
    \item<2-> R
    \begin{itemize}
        \item De facto standard in statistics
        \item Great for work with real data
        \item Matrix syntax is less intuitive
    \end{itemize}
    \item<3-> Python
    \begin{itemize}
        \item De facto standard in the tech industry and, increasingly, physical sciences
        \item Great all-purpose language (``Swiss army knife'')
        \item Matrix syntax has improved but is still less intuitive than Julia's
    \end{itemize}
\end{itemize}
\end{frame}

\begin{frame}{Pros and cons of Julia}
\begin{itemize}
    \item Pros\itemsep=2em
    \begin{itemize}
        \item Clean syntax and fast execution for numerical computation\itemsep=2em
        \item Native support for automatic differentiation makes optimization easy, robust, and quick
    \end{itemize}
    \item<2-> Cons
    \begin{itemize}
        \item Generally harder to use than R or Python for manipulation of real data\itemsep=2em
        \item The language itself changes regularly (most resources from before 2018/Julia v1.0 are useless)
        \item The community is smaller than R and Python
    \end{itemize}
\end{itemize}
\end{frame}

\begin{frame}{More on Julia vs. other languages}
\begin{itemize}
    \item \href{https://tradediversion.net/2018/09/17/why-i-encourage-econ-phd-students-to-learn-julia/}{Why I encourage econ PhD students to learn Julia} (Jonathan Dingel, September 2018)\itemsep=2em
    \item \href{https://www.sas.upenn.edu/~jesusfv/Lecture_HPC_5_Scientific_Computing_Languages.pdf}{Scientific Computing Languages} (Jesús Fernández-Villaverde, November 2019)
\end{itemize}
\end{frame}

\begin{frame}{Now...}
Time for the demo!
\end{frame}
\end{document}